\section{Appendix B. Implementation Details}

\subsection{B.1 Core Module}

\begin{lstlisting}
class HGRN2Layer(BaseAttention):
    def \_\_init\_\_(self, config, rank=None, d\_model=None, 
                 res\_clip=30.0, lora\_rank=16, lora\_alpha=32.0, 
                 is\_decoder=None, freeze\_base=True, **kwargs):
        super().\_\_init\_\_(config, rank)
        
        self.d\_model = resolve\_d\_model(config, d\_model, kwargs)
        self.res\_clip = res\_clip
        
        # Gates
        self.forget\_gate = nn.Linear(self.d\_model, self.d\_model)
        self.input\_gate = nn.Linear(self.d\_model, self.d\_model)
        self.v\_proj = nn.Linear(self.d\_model, self.d\_model, bias=False)
        self.g\_proj = nn.Linear(self.d\_model, self.d\_model)
        self.state\_norm = nn.LayerNorm(self.d\_model)
        self.out\_proj = nn.Linear(self.d\_model, self.d\_model)
        
        # Init
        nn.init.normal\_(self.forget\_gate.weight, std=0.01)
        nn.init.constant\_(self.forget\_gate.bias, -2.5)
        nn.init.normal\_(self.input\_gate.weight, std=0.02)
        nn.init.constant\_(self.input\_gate.bias, 0.0)
        nn.init.normal\_(self.g\_proj.weight, std=0.02)
        nn.init.constant\_(self.g\_proj.bias, 1.28)
        nn.init.normal\_(self.v\_proj.weight, std=0.02)
\end{lstlisting}

\subsection{B.2 The Chunked Scan}

\begin{lstlisting}
def \_scan\_direction(self, f, iv, chunk\_size=32):
    B, L, D = iv.shape
    f\_fp32 = f.float()
    iv\_fp32 = iv.float()
    log\_f = torch.log(f\_fp32.clamp(min=1e-6))
    
    h\_chunks = []
    h\_prev = torch.zeros(B, D, device=iv.device, dtype=torch.float32)
    
    for start in range(0, L, chunk\_size):
        end = min(start + chunk\_size, L)
        C = end - start
        
        f\_c = log\_f[:, start:end]           # (B, C, D)
        iv\_c = iv\_fp32[:, start:end]        # (B, C, D)
        
        A = torch.cumsum(f\_c, dim=1)        # (B, C, D)
        A\_diff = A.unsqueeze(2) - A.unsqueeze(1)   # (B, C, C, D)
        
        mask = torch.tril(torch.ones(C, C, device=iv.device, dtype=torch.bool))
        A\_diff = torch.where(mask.unsqueeze(0).unsqueeze(-1), A\_diff, 
                             torch.full\_like(A\_diff, -1e4))
        
        decay\_weights = torch.exp(A\_diff) * mask.unsqueeze(0).unsqueeze(-1)
        intra = torch.sum(decay\_weights * iv\_c.unsqueeze(1), dim=2)
        inter = h\_prev.unsqueeze(1) * torch.exp(A)
        h\_c = inter + intra
        
        if self.res\_clip is not None and self.res\_clip \textgreater{} 0:
            h\_c = torch.clamp(h\_c, -self.res\_clip, self.res\_clip)
        
        h\_chunks.append(h\_c)
        h\_prev = h\_c[:, -1]
    
    return torch.cat(h\_chunks, dim=1).to(iv.dtype)
\end{lstlisting}

\subsection{B.3 Forget Gate Initialization}

\begin{lstlisting}
nn.init.normal\_(self.forget\_gate.weight, mean=0.0, std=0.01)
nn.init.constant\_(self.forget\_gate.bias, -2.5)
\end{lstlisting}

\textbf{Why -2.5?} \texttt{softplus(-2.5) $\approx$ 0.076}. So \texttt{log f = -0.076}, 
\texttt{f = e^{-0.076} $\approx$ 0.93}. This means the state retains 93\% of its value 
per step -- slow decay.

\subsection{B.4 Input Modulation}

\begin{lstlisting}
input\_scale = (1.0 - f) * i
iv = input\_scale * v
\end{lstlisting}

*\textit{Why (1-f) } i?** When $f \approx 1$ (retain), $1 - f \approx 0$, so no 
input. When $f \approx 0$ (forget), $1 - f \approx 1$, so full input. This 
ties the write magnitude to the forget rate.

\subsection{B.5 Output Modulation}

\begin{lstlisting}
h\_norm = self.state\_norm(h\_ctx)
gated = (0.5 * v + 0.5 * h\_norm) * F.silu(z)
out = self.out\_proj(gated)
\end{lstlisting}

\textbf{Why 0.5 v + 0.5 h\_norm?} Preserves the current value alongside the 
recurrent state.

\subsection{B.6 Failure Modes}

\subsubsection{B.6.1 Missing V/O}

If V/O weights cannot be extracted, \texttt{copy\_weights} logs a warning but does 
not raise. This is a known limitation.

\subsubsection{B.6.2 Nested Encoder}

For nested encoders, we set \texttt{layer.attention.output.dense = nn.Identity()}. 
If this is wrong, the output projection is skipped.

\subsubsection{B.6.3 Cross-Attention}

Cross-attention raises \texttt{NotImplementedError}. This is intentional.

\subsection{B.7 Common Pitfalls}

\subsubsection{B.7.1 Random Gate Initialization}

Without careful init, forget gates can saturate.

\subsubsection{B.7.2 FP16 Scan}

FP16 scan loses precision. We use FP32 for the scan.

\subsubsection{B.7.3 Missing Clamp}

Without \texttt{res\_clip}, states can grow unbounded.

\subsubsection{B.7.4 Wrong Nested Handling}

For RoBERTa, forgetting to set \texttt{output.dense = nn.Identity()} causes 
double projection.
\begin{lstlisting}

---
\end{lstlisting}